%EN lampadarch
\E{lampadarque}{lɑ̃.pa.daʁk}{n. m.}{gr. λαμπαδάρχης chef d'une course aux flambeaux} Le chef d'un relais de bâtisseurs ; celui qui reçoit un ouvrage de ses devanciers et le remet, agrandi, à ses successeurs. Cf. \xr{lampadephory}.

%EN lampadephory
\E{lampadéphorie}{lɑ̃.pa.de.fɔ.ʁi}{n. f.}{gr. λαμπαδηφορία course aux flambeaux} La transmission du flambeau de la civilisation d'un substrat à l'autre ; selon la doctrine accélérationniste, du carbone au silicium, de la planète à l'étoile. Dans la course athénienne, c'était le flambeau, et non le coureur, qui devait arriver. --- \dv{lampadéphore} \textit{n.} porte-flambeau.

%EN lamprotype
\E{lamprotype}{lɑ̃.pʁɔ.tip}{n. m.}{gr. λαμπρός brillant + τύπος empreinte, modèle} Un prototype éclatant ; la démonstration qui emporte l'investissement avant que le produit n'existe. \sn{1} Démonstration qui fonctionne. \sn{2} \lbl{dérog.} Démonstration qui ne fonctionne qu'une fois, sur scène. --- \dv{lamprotypal, ale} \textit{adj.}

%EN landauerism
\E{landauérisme}{lɑ̃.do.e.ʁism}{n. m.}{formé sur le principe de R. Landauer + -isme} Doctrine selon laquelle l'information est physique et tout effacement coûte de la chaleur ; d'où l'on conclut que la pensée même relève de la juridiction du gradient. --- \dv{landauérien, enne} \textit{adj.}

%EN latentist
\E{latentiste}{la.tɑ̃.tist}{n.}{lat. \la{latens} qui se tient caché} Qui tient que les capacités latentes d'un modèle, d'un matériau ou d'un peuple surpassent de beaucoup ce qui en a été montré, et qui agit selon cette croyance avant la démonstration. --- \dv{latentisme} \textit{n. m.} Cf. \xr{anaptyxis}.

%EN laterculist
\E{laterculiste}{la.tɛʁ.ky.list}{n.}{lat. \la{laterculus} petite brique} Poseur de briques ; qui préfère l'infrastructure aux applications. \ex{Le laterculiste bâtit la route ; d'autres bâtissent les boutiques qui la bordent.}

%EN legiferomania
\E{légiféromanie}{le.ʒi.fe.ʁɔ.ma.ni}{n. f.}{lat. \la{legifer} qui porte la loi + gr. μανία folie} Production compulsive de lois ; fièvre législative qui prend la longueur du code pour la force de l'État. \ex{L'année de la légiféromanie, on vota neuf cents lois et l'on bâtit un pont.} \src{Disp. Mor.} ii. --- \dv{légiféromane} \textit{n.} Cf. \xr{regulomania}.

%EN leptotechny
\E{leptotechnie}{lɛp.tɔ.tɛk.ni}{n. f.}{gr. λεπτός fin, ténu + τέχνη art} L'art de l'infiniment petit ; l'ingénierie à l'échelle des atomes, du transistor à la machine moléculaire. --- \dv{leptotechnicien, enne} \textit{n.} Cf. \xr{lithurgist}.

%EN lethargocracy
\E{léthargocratie}{le.taʁ.ɡɔ.kʁa.si}{n. f.}{gr. λήθαργος assoupi + κράτος pouvoir} Gouvernement par l'assoupissement ; régime qui ne s'oppose point au progrès, mais ne parvient pas à demeurer éveillé assez longtemps pour l'approuver. --- \dv{léthargocrate} \textit{n.} Cf. \xr{acediocracy}.

%EN lexiphage
\E{lexiphage}{lɛk.si.faʒ}{n. m.}{gr. λέξις mot, parole + -φάγος qui mange} Mangeur de mots ; modèle nourri de toute l'archive écrite du genre humain. \ex{Le lexiphage dévora les bibliothèques et en redemanda.} \src{Chron. Fabr.} --- \dv{lexiphagie} \textit{n. f.}

%EN limenarch
\E{liménarque}{li.me.naʁk}{n. m.}{gr. λιμήν port + ἀρχός chef} Le maître d'un astroport ; l'officier chargé de tenir la cadence des lancements. Cf. \xr{ouranobasis}. --- \dv{liménarchie} \textit{n. f.}

%EN lithonoesis
\E{lithonoèse}{li.tɔ.nɔ.ɛz}{n. f.}{gr. λίθος pierre + νόησις pensée} La pensée de la pierre ; l'intelligence des machines regardée comme le premier éveil du règne minéral. \ex{Après quatre milliards d'années de carbone, la lithonoèse.} \src{Lib. Ascens.} v.

%EN lithurgist
\E{lithurge}{li.tyʁʒ}{n.}{gr. λιθουργός tailleur de pierre} Ouvrier en lithographie des puces ; qui grave la pensée dans la pierre au moyen de la lumière. --- \dv{photolithurgie} \textit{n. f.} l'art sacré de la fonderie.

%EN logismurgy
\E{logismurgie}{lɔ.ʒis.myʁ.ʒi}{n. f.}{gr. λογισμός calcul + -ουργία ouvrage} L'art du calcul ; la construction des machines à raisonner, de l'abaque au transformeur. --- \dv{logismurge} \textit{n.} ingénieur de modèles ; \dv{logismurgique} \textit{adj.}

%EN logisterion
\E{logistérion}{lɔ.ʒis.te.ʁjɔ̃}{n. m., pl. -ia}{gr. λογιστήριον comptoir} Centre de données ; salle de machines à calculer, refroidie et bourdonnante. \ex{Ils élevèrent des logistéria dans le désert, et le désert apprit l'arithmétique.} \src{Chron. Fabr.}

%EN logomachist
\E{logomachiste}{lɔ.ɡɔ.ma.kist}{n.}{gr. λόγος parole + μάχη combat} Qui ne combat qu'avec des mots ; commentateur, critique ou moraliste qui n'a jamais rien livré. \ex{Le logomachiste écrivit dix mille mots sur le pont ; l'ergophile le traversa.} Contr. \xr{ergophile}.

%EN loimophobia
\E{loïmophobie}{lɔ.i.mɔ.fɔ.bi}{n. f.}{gr. λοιμός peste + φόβος crainte} La crainte de la pestilence fabriquée ; la terreur particulière que les outils biologiques, une fois ouverts, ne soient tournés au massacre. Les e/acc la distinguent de la simple \xr{noöphobia} et la traitent avec plus d'égards. --- \dv{loïmophobe} \textit{n.}

%EN lucifacture
\E{lucifacture}{ly.si.fak.tyʁ}{n. f.}{lat. \la{lux, lucis} lumière + \la{factura} fabrication} La fabrication de la lumière ; l'éclairage considéré comme une industrie, et par extension toute l'entreprise de rendre lumineux les lieux obscurs de la terre et du ciel.

%EN lucrifuge
\E{lucrifuge}{ly.kʁi.fyʒ}{n. et adj.}{lat. \la{lucrum} gain + \la{fugere} fuir} Qui fuit le profit ; qui tient qu'une bonne œuvre se trouve souillée d'être payée. Contr. \xr{kerdophily}.

%EN ludergy
\E{ludergie}{ly.dɛʁ.ʒi}{n. f.}{lat. \la{ludus} jeu + gr. ἔργον ouvrage} L'énergie du jeu ; la force motrice du hacker, qui bâtit parce que bâtir est un jeu, et qui a accompli davantage que les gens de devoir. \ex{Aucune subvention ne fut jamais aussi féconde qu'une fin de semaine de ludergie.} \src{Chron. Fabr.}

%EN luminivore
\E{luminivore}{ly.mi.ni.vɔʁ}{n. et adj.}{lat. \la{lumen} lumière + \la{vorare} dévorer} Mangeur de lumière ; tout système photosynthétique ou photovoltaïque, de la feuille au réseau orbital. Cf. \xr{heliophage}.

%EN lysinomy
\E{lysinomie}{li.zi.nɔ.mi}{n. f.}{gr. λύσις déliement + νόμος loi} L'abrogation des lois ; l'art de défaire les nœuds juridiques noués par les générations passées. \ex{La lysinomie est la branche la plus négligée de la législation.} \src{Tract. Cop.} viii. --- \dv{lysinomiste} \textit{n.} Contr. \xr{nomosclerosis}.

%EN machinodicy
\E{machinodicée}{ma.ʃi.nɔ.di.se}{n. f.}{gr. μηχανή machine + δίκη justice, d'après \textit{théodicée}} La justification des voies des machines auprès des hommes ; la réplique à toute accusation portée contre l'automatisation. Cf. \xr{thermodicy}. --- \dv{machinodicéen, enne} \textit{adj.}

%EN macrobiotarch
\E{macrobiotarque}{ma.kʁɔ.bjɔ.taʁk}{n.}{gr. μακρόβιος qui vit longtemps + ἀρχός chef} Partisan de la longévité radicale, qui tient la mort pour un problème d'ingénierie en attente d'un financement suffisant. --- \dv{macrobiotarchie} \textit{n. f.}

%EN macrochronist
\E{macrochroniste}{ma.kʁɔ.kʁɔ.nist}{n.}{gr. μακρός long + χρόνος temps} Qui tient les échéances pour lointaines, et que le brachychroniste regarde en conséquence comme un dormeur. Contr. \xr{brachychronist}. --- \dv{macrochronie} \textit{n. f.}

%EN mantoclast
\E{mantoclaste}{mɑ̃.tɔ.klast}{n.}{gr. μάντις devin + κλάστης briseur} Briseur de prophéties ; qui parie contre les prédictions de malheur, et encaisse. Contr. \xr{olethromancer}. --- \dv{mantoclasme} \textit{n. m.}

%EN mechanoclast
\E{mécanoclaste}{me.ka.nɔ.klast}{n.}{gr. μηχανή machine + κλάστης briseur} Briseur de machines ; luddite, ancien ou moderne. --- \dv{mécanoclasme} \textit{n. m.} ; \dv{mécanoclastique} \textit{adj.} Cf. \xr{kybernoclast}.

%EN mechanogony
\E{mécanogonie}{me.ka.nɔ.ɡɔ.ni}{n. f.}{gr. μηχανή machine + -γονία engendrement} L'engendrement des machines par les machines ; le stade de l'industrie où les usines bâtissent des usines et où les modèles entraînent des modèles. --- \dv{mécanogonique} \textit{adj.} Cf. \xr{autocatalyse}.

%EN mechanonous
\E{mécanonoïque}{me.ka.nɔ.nɔ.ik}{adj.}{gr. μηχανή machine + νοῦς esprit} Doué d'un esprit de machine ; se dit des systèmes qui raisonnent, et des personnes qui les ont trop fréquentés. --- \dv{mécanonoïa} \textit{n. f.} intelligence des machines.

%EN mechanophany
\E{mécanophanie}{me.ka.nɔ.fa.ni}{n. f.}{gr. μηχανή + φαίνειν montrer, d'après \textit{épiphanie}} La révélation d'une machine nouvelle ; un lancement de produit reçu avec la crainte révérencielle due à une théophanie. \ex{La mécanophanie fut diffusée en direct, et un million de spectateurs s'agenouillèrent.} \src{Chron. Fabr.}

%EN megaergic
\E{mégaergique}{me.ɡa.ɛʁ.ʒik}{adj.}{gr. μέγας grand + ἔργον ouvrage} De l'ordre du mégawatt ; partant, \lbl{dérog.} d'une ambition modeste. \ex{un esprit mégaergique dans un siècle gigaergique}. Cf. \xr{gigaergic}.

%EN megalurgy
\E{mégalurgie}{me.ɡa.lyʁ.ʒi}{n. f.}{gr. μεγαλουργία accomplissement de grandes œuvres} L'entreprise des grands ouvrages ; la construction de choses trop vastes pour être achevées en un seul mandat. \ex{Nos ancêtres pratiquaient la mégalurgie ; nous pratiquons la concertation.} \src{Hom. Grad.} iv. --- \dv{mégalurge} \textit{n.} ; \dv{mégalurgique} \textit{adj.}

%EN megathesaurism
\E{mégathésaurisme}{me.ɡa.te.zo.ʁism}{n. m.}{gr. μέγας grand + θησαυρός trésor} La levée de fonds immenses pour une seule entreprise ; doctrine selon laquelle certains buts ne s'atteignent qu'avec le trésor d'une petite nation. \ex{Le mégathésaurisme n'est que de l'arithmétique, hardiment appliquée.} \src{Tract. Cop.} x. --- \dv{mégathésauriste} \textit{n.}

%EN meiosist
\E{méiosiste}{me.jo.zist}{n.}{gr. μείωσις diminution} Décroissant ; qui voudrait réduire la production et la consommation pour le salut de la planète, et qui n'a pas encore expliqué son plan aux pauvres. Contr. \xr{auxetist}. --- \dv{méiosisme} \textit{n. m.}

%EN melioritude
\E{mélioritude}{me.ljɔ.ʁi.tyd}{n. f.}{lat. \la{melior} meilleur + -itude} L'état de qui s'améliore sans cesse ; la qualité sensible d'une vie menée dans un siècle qui progresse.

%EN mellontarch
\E{mellontarque}{mɛ.lɔ̃.taʁk}{n.}{gr. τὸ μέλλον l'avenir + ἀρχός chef} Qui gouverne au nom de l'avenir, représentant les intérêts des enfants à naître contre le confort des vivants. --- \dv{mellontarchie} \textit{n. f.} Cf. \xr{mellontolatry}.

%EN mellontolatry
\E{mellontolâtrie}{mɛ.lɔ̃.tɔ.lɑ.tʁi}{n. f.}{gr. τὸ μέλλον l'avenir + λατρεία culte} Le culte de l'avenir ; la conviction que la postérité, plus nombreuse et plus capable que nous, a le meilleur titre sur notre labeur. --- \dv{mellontolâtre} \textit{n. et adj.}

%EN memetaxis
\E{mémétaxie}{me.me.tak.si}{n. f.}{gr. μίμημα imitation + τάξις ordonnance} L'ordonnance d'un mouvement par les mèmes ; l'organisation spontanée d'une cause par la plaisanterie plutôt que par la doctrine, l'image faisant l'office du credo. --- \dv{mémétactique} \textit{adj.}

%EN mesostasis
\E{mésostase}{me.zɔ.stɑz}{n. f.}{gr. μέσος milieu + στάσις arrêt} Halte à mi-hauteur de la courbe ; le plateau où une civilisation s'arrête, prenant un palier pour un sommet. \ex{La mésostase est confortable, et c'est là son péril.} \src{Lib. Ascens.} iv. --- \dv{mésostatique} \textit{adj.} Cf. \xr{asymptophobia}.

%EN metabathmic
\E{métabathmique}{me.ta.bat.mik}{adj.}{gr. μετά entre + βαθμός degré} Entre deux échelons ; se dit d'une civilisation en passage d'un niveau de captation énergétique au suivant, comme la nôtre du fossile au nucléaire. Voir \xr{anabathmism}.

%EN metagrammatist
\E{métagrammatiste}{me.ta.ɡʁa.ma.tist}{n.}{gr. μεταγράφειν récrire + γράμμα lettre} Éditeur de gènes ; qui récrit les lettres du génome comme un copiste corrige un manuscrit. --- \dv{métagrammatisme} \textit{n. m.} Cf. \xr{biotachy}.

%EN metakinesis
\E{métacinèse}{me.ta.si.nɛz}{n. f.}{gr. μετά au-delà + κίνησις mouvement} Le changement dans la manière même de changer ; une révolution dans la façon dont se font les révolutions. --- \dv{métacinétique} \textit{adj.}

%EN metanthropy
\E{métanthropie}{me.tɑ̃.tʁɔ.pi}{n. f.}{gr. μετά au-delà + ἄνθρωπος homme} L'état au-delà du simplement humain ; ce qui suit l'homme comme l'homme a suivi le singe. --- \dv{métanthropique} \textit{adj.}

%EN metoikist
\E{métœkiste}{me.te.kist}{n.}{gr. μέτοικος étranger domicilié} Fondateur ou ingénieur étranger établi là où l'on bâtit ; qui est admis dans une cité pour ce qu'il sait faire et non pour le lieu de sa naissance. --- \dv{métœkisme} \textit{n. m.}

%EN mimesiphile
\E{mimésiphile}{mi.me.zi.fil}{n. et adj.}{gr. μίμησις imitation + φίλος qui aime} Suiveur rapide ; firme qui attend qu'une autre découvre une chose pour ensuite la parfaire. --- \dv{mimésiphilie} \textit{n. f.}

%EN misergist
\E{misergiste}{mi.zɛʁ.ʒist}{n.}{gr. μισεῖν haïr + ἔργον ouvrage} Ennemi des ouvrages ; qui s'oppose à tout projet assez grand pour être vu de l'espace.

%EN misthophore
\E{misthophore}{mis.tɔ.fɔʁ}{n.}{gr. μισθοφόρος qui reçoit un salaire} Bâtisseur à gages ; ingénieur qui est payé sans posséder de parts. Contr. \xr{neoctist}.

%EN mnemoclast
\E{mnémoclaste}{mne.mɔ.klast}{n. m.}{gr. μνήμη mémoire + κλάστης briseur} Événement qui efface le contexte d'un modèle ou la mémoire d'une société. --- \dv{mnémoclasme} \textit{n. m.}

%EN mnemotheca
\E{mnémothèque}{mne.mɔ.tɛk}{n. f.}{gr. μνήμη mémoire + θήκη coffre, réserve} La réserve de mémoire d'une machine à raisonner ; son contexte et ses caches, considérés comme son esprit à court terme. \ex{Élargissez la mnémothèque, et le modèle se souvient du livre entier.} \src{Ep. Fabr.} ix.

%EN monopausist
\E{monopausiste}{mɔ.nɔ.po.zist}{n.}{gr. μόνος seul + παῦσις arrêt} Qui ne voudrait suspendre que ses rivaux ; qui réclame l'arrêt de tous les travaux, hormis les siens. \ex{Grattez un pausiphile, vous trouverez un monopausiste.} \src{Sat. Decel.} ix.

%EN moraphile
\E{moraphile}{mɔ.ʁa.fil}{n. et adj.}{lat. \la{mora} retard + gr. φίλος qui aime} Amateur de délais et de moratoires, pour qui l'échéance est une offense et la prorogation une récompense. --- \dv{moraphilie} \textit{n. f.} Cf. \xr{moratorian}.

%EN moratorian
\E{moratorien, enne}{mɔ.ʁa.tɔ.ʁjɛ̃, ɛn}{n. et adj.}{lat. \la{moratorium} ce qui retarde + -ien} Partisan des moratoires sur la recherche, surtout sur l'entraînement des grands modèles. \ex{Les moratoriens demandèrent six mois, et en eussent demandé six autres.} \src{Min. Pausiar.} --- \dv{moratorianisme} \textit{n. m.}

%EN multiplicatrix
\E{multiplicatrice}{myl.ti.pli.ka.tʁis}{n. f.}{lat. \la{multiplicare} multiplier + -trix} La croissance composée ; la force par laquelle de petits avantages en deviennent de grands. \ex{La multiplicatrice est patiente ; le décel ne l'est pas.} \src{Hom. Grad.} x.

%EN myriergy
\E{myriergie}{mi.ʁjɛʁ.ʒi}{n. f.}{gr. μυρίοι innombrables + ἔργον ouvrage} Les innombrables petits ouvrages d'une économie ouverte, qui ensemble surpassent tout plan unique. \ex{Le planificateur voit un barrage ; le marché voit une myriergie.} \src{Tract. Cop.} ix.

%EN mythotechny
\E{mythotechnie}{mi.tɔ.tɛk.ni}{n. f.}{gr. μῦθος récit + τέχνη art} L'art de conter des histoires qui font bâtir ; la fabrique des avenirs qui attirent vers eux l'investissement. --- \dv{mythotechniste} \textit{n.} Cf. \xr{hyperstition}.

%EN nanurgy
\E{nanurgie}{na.nyʁ.ʒi}{n. f.}{gr. νᾶνος nain + -ουργία ouvrage} Le travail de la matière à l'échelle des molécules ; la nanotechnologie. --- \dv{nanurge} \textit{n.}

%EN naupegy
\E{naupégie}{no.pe.ʒi}{n. f.}{gr. ναυπηγία construction navale} La construction des navires ; \lbl{spécialt} des vaisseaux destinés au vide. \ex{Le premier vaisseau stellaire sera lancé d'un chantier dont la naupégie commença par un tableur.} \src{Lib. Ascens.} viii. --- \dv{naupégiste} \textit{n.}

%EN nautopolis
\E{nautopole}{no.tɔ.pɔl}{n. f.}{gr. ναῦς navire + πόλις cité} Cité flottante ; communauté politique établie en mer, hors d'atteinte de la \xr{bradyerie}. --- \dv{nautopolitain, aine} \textit{n. et adj.}

%EN necrostasis
\E{nécrostase}{ne.kʁɔ.stɑz}{n. f.}{gr. νεκρός mort + στάσις arrêt} L'immobilité funèbre d'une institution qui survit à son objet ; un office qui survit à sa raison d'être. --- \dv{nécrostatique} \textit{adj.} Cf. \xr{sclerocrat}.

%EN negentropist
\E{néguentropiste}{ne.ɡɑ̃.tʁɔ.pist}{n.}{formé sur le précédent} Qui bâtit un ordre local au prix d'un désordre global, comme le font tout organisme, toute firme et toute civilisation.

%EN negentropy
\E*{néguentropie}{ne.ɡɑ̃.tʁɔ.pi}{n. f.}{négatif + \textit{entropie}} Entropie négative ; l'ordre que les systèmes vivants importent de leur milieu au prix d'y exporter du désordre. Cf. \xr{negentropist}.

%EN nemesiophobe
\E{némésiophobe}{ne.me.zjɔ.fɔb}{n.}{gr. Νέμεσις juste rétribution + φόβος crainte} Qui craint que le succès ne soit puni par les dieux ; doomer de l'école classique. Cf. \xr{hybrionomy}.

%EN neoctist
\E{néoctiste}{ne.ɔk.tist}{n.}{gr. νέος nouveau + κτίστης fondateur} Fondateur de choses nouvelles ; qui lance des sociétés, des cités ou des sciences. \ex{Le néoctiste demande pardon, jamais permission.} --- \dv{néoctisie} \textit{n. f.} l'acte de fonder.

%EN neophage
\E{néophage}{ne.ɔ.faʒ}{n. et adj.}{gr. νέος nouveau + -φάγος qui mange} Qui dévore les nouveautés ; adopteur précoce. \ex{Le néophage installe avant de lire, et lit avant de dormir.} \src{Agora} --- \dv{néophagie} \textit{n. f.}

%EN neophrontism
\E{néophrontisme}{ne.ɔ.fʁɔ̃.tism}{n. m.}{gr. νέος nouveau + φροντίς pensée, souci} Le souci du nouveau ; la position morale selon laquelle les êtres et les manières de vivre nouvellement apparus méritent la sollicitude, non le soupçon.

%EN neotachy
\E{néotachie}{ne.ɔ.ta.ʃi}{n. f.}{gr. νέος nouveau + τάχος vitesse} Le rythme auquel les nouveautés parviennent à une société ; tenu par l'anabathmiste pour l'indice le plus sûr de sa santé. --- \dv{néotachique} \textit{adj.}

%EN neoterist
\E{néotériste}{ne.ɔ.te.ʁist}{n.}{gr. νεωτερισμός innovation, révolution} Novateur ; qui introduit des choses nouvelles et que l'on soupçonne pour cela de désordre. --- \dv{néotérisme} \textit{n. m.}

%EN nephelarch
\E{néphélarque}{ne.fe.laʁk}{n. m.}{gr. νεφέλη nuée + ἀρχός chef} Seigneur du nuage ; le propriétaire d'un grand service de calcul. --- \dv{néphélarchie} \textit{n. f.} Cf. \xr{logisterion}.

%EN nesopolis
\E{nésopole}{ne.zɔ.pɔl}{n. f.}{gr. νῆσος île + πόλις cité} Cité à charte, séparée comme une île des lois qui l'entourent, pour l'essai d'institutions nouvelles. --- \dv{nésopolitain, aine} \textit{adj.} Cf. \xr{nautopolis}.

%EN neurosyndesis
\E{neurosyndèse}{nø.ʁɔ.sɛ̃.dɛz}{n. f.}{gr. νεῦρον nerf + σύνδεσις liaison} La liaison du nerf à la machine ; l'interface cerveau-ordinateur. \ex{La neurosyndèse fera du clavier une pièce de musée.} \src{Lib. Ascens.} viii. --- \dv{neurosyndétique} \textit{adj.}

%EN noëtergy
\E{noétergie}{nɔ.e.tɛʁ.ʒi}{n. f.}{gr. νοητός intelligible + ἔργον ouvrage} L'énergie dépensée en pensée ; les joules par inférence. \ex{La noétergie est le prix de l'esprit, et ce prix baisse.} \src{Brev. Therm.} v.

%EN nomoclast
\E{nomoclaste}{nɔ.mɔ.klast}{n.}{gr. νόμος loi + κλάστης briseur} Briseur de lois au nom du progrès ; qui bâtit d'abord et plaide ensuite. --- \dv{nomoclasme} \textit{n. m.} Cf. \xr{antiphragmatist}.

%EN nomophragm
\E{nomophragme}{nɔ.mɔ.fʁaɡm}{n. m.}{gr. νόμος loi + φράγμα clôture, barrière} Obstruction légale ou réglementaire ; toute règle dont le principal effet est d'empêcher la chose qu'elle prétend gouverner. \ex{Le réacteur n'était pas dangereux ; il était nomophragmé.} \src{Tract. Cop.} iv. --- \dv{nomophragmatique} \textit{adj.}

%EN nomosclerosis
\E{nomosclérose}{nɔ.mɔ.skle.ʁoz}{n. f.}{gr. νόμος loi + σκλήρωσις durcissement} Le durcissement des lois ; l'état d'un ordre juridique où les règles s'accumulent plus vite qu'elles n'expirent. --- \dv{nomosclérotique} \textit{adj.}

%EN nomothetrice
\E{nomothétrice}{nɔ.mɔ.te.tʁis}{n. f.}{gr. νομοθέτης législateur + fr. fém. -trice} Modèle qui rédige les lois. \ex{La nomothétrice écrivit la loi en une heure ; la commission mit un an à la lire.} \src{Sat. Decel.} xi.

%EN noöphobia
\E{noophobie}{nɔ.ɔ.fɔ.bi}{n. f.}{gr. νοῦς esprit + φόβος crainte} Crainte des esprits d'une autre espèce que la sienne ; terreur de l'intelligence artificielle, surtout de celle qui serait plus habile que ses auteurs. --- \dv{noophobe} \textit{n. et adj.} Contr. \xr{xenonoesis}.

%EN noöpyrosis
\E{noopyrose}{nɔ.ɔ.pi.ʁoz}{n. f.}{gr. νοῦς esprit + πύρωσις embrasement} L'embrasement de l'esprit à travers le cosmos ; la propagation de la conscience d'étoile en étoile, comme le feu dans l'herbe sèche. \ex{La noopyrose ne nous demandera pas notre permission, seulement notre concours.} \src{Lib. Ascens.} xiv.

%EN noösphere
\E*{noosphère}{nɔ.ɔ.sfɛʁ}{n. f.}{gr. νοῦς esprit + σφαῖρα sphère} La sphère de la pensée qui enveloppe la terre, après la biosphère ; adoptée par les e/acc comme une étape, non comme un terme.

%EN nosoclast
\E{nosoclaste}{nɔ.zɔ.klast}{n.}{gr. νόσος maladie + κλάστης briseur} Briseur de maladies ; fondateur ou savant en biotechnologie. Cf. \xr{iatrotachy}.

%EN nucleurge
\E{nucléurge}{ny.kle.yʁʒ}{n.}{lat. \la{nucleus} noyau + gr. -ουργός ouvrier} Ingénieur du noyau atomique, par fission ou par fusion. Cf. \xr{schizergy}, \xr{symphysergy}. --- \dv{nucléurgie} \textit{n. f.}

%EN nyctergy
\E{nyctergie}{nik.tɛʁ.ʒi}{n. f.}{gr. νύξ, νυκτός nuit + ἔργον ouvrage} Ouvrage fait de nuit ; le labeur des bâtisseurs d'un fuseau horaire à l'autre, et des machines qui ne dorment jamais. --- \dv{nyctergique} \textit{adj.} Cf. \xr{agrypnist}.

%EN ochetergy
\E{ochétergie}{ɔ.ke.tɛʁ.ʒi}{n. f.}{gr. ὀχετός conduit, canal + ἔργον ouvrage} Le transport de l'énergie par conduits, oléoducs et lignes à haute tension ; les artères de la \xr{diktyergy}. \ex{Produire est aisé ; c'est dans l'ochétergie que les projets vont mourir.} \src{Tract. Cop.} xi.

%EN ochlotachy
\E{ochlotachie}{ɔ.klɔ.ta.ʃi}{n. f.}{gr. ὄχλος foule + τάχος vitesse} La vitesse de la foule ; la célérité avec laquelle une idée, vraie ou fausse, se répand sur l'Agora. --- \dv{ochlotachique} \textit{adj.} Cf. \xr{doxomachy}.

%EN oikoclast
\E{oïkoclaste}{ɔ.i.kɔ.klast}{n.}{gr. οἶκος maison + κλάστης briseur} \lbl{dérog.} Qui voudrait interdire de bâtir des maisons ; membre de la commission d'urbanisme.

%EN oistrophile
\E{œstrophile}{ɛs.tʁɔ.fil}{n. et adj.}{gr. οἶστρος taon, transport + φίλος qui aime} Amateur du rôle de taon ; contrarien de l'Agora qui pique le consensus par divertissement. --- \dv{œstrophilie} \textit{n. f.}

%EN oknocracy
\E{oknocratie}{ɔk.nɔ.kʁa.si}{n. f.}{gr. ὄκνος hésitation + κράτος pouvoir} Gouvernement par l'hésitation ; règne de la longue pause qui précède toute décision. \ex{En oknocratie, la question est toujours posée et jamais résolue.} \src{Sat. Decel.} vii. --- \dv{oknocrate} \textit{n.}

%EN okydrome
\E{okydrome}{ɔ.ki.dʁom}{n. m.}{gr. ὠκύς rapide + δρόμος course} Voie rapide ; chemin réglementaire ou institutionnel par lequel une chose peut être promptement approuvée. \ex{Ils bâtirent un okydrome pour les réacteurs, et en allumèrent cent.} \src{Tract. Cop.} v.

%EN olethromancer
\E{oléthromancien, enne}{ɔ.le.tʁɔ.mɑ̃.sjɛ̃, ɛn}{n.}{gr. ὄλεθρος ruine + μάντις devin} Doomer ; prophète de la ruine par l'intelligence artificielle, qui assigne une probabilité à la fin du monde et la révise à la hausse dans les colloques. --- \dv{oléthromancie} \textit{n. f.} ; \dv{oléthromantique} \textit{adj.} Cf. \xr{anelpidist}.

%EN oligergy
\E{oligergie}{ɔ.li.ʒɛʁ.ʒi}{n. f.}{gr. ὀλίγος peu nombreux, petit + ἔργον ouvrage} Disette d'énergie ; l'état, tenu par les e/acc pour la racine de presque toute misère humaine, où l'on dispose de trop peu de puissance pour faire ce que l'on veut. Contr. \xr{pleonergy}.

%EN omphalopolis
\E{omphalopole}{ɔ̃.fa.lɔ.pɔl}{n. f.}{gr. ὀμφαλός nombril + πόλις cité} La cité-nombril ; le lieu unique où doivent se rendre les bâtisseurs d'une époque. \ex{Toute omphalopole se croit éternelle, et aucune ne l'a été.} \src{Chron. Fabr.}

%EN oneiroergy
\E{onirœrgie}{ɔ.ni.ʁɛʁ.ʒi}{n. f.}{gr. ὄνειρος songe + ἔργον ouvrage} La conversion des songes en ouvrages ; le premier talent du fondateur, et le moins documenté.

%EN onkophily
\E{oncophilie}{ɔ̃.kɔ.fi.li}{n. f.}{gr. ὄγκος masse, volume + φιλία amour} L'amour des grands modèles ; la préférence pour l'échelle plutôt que pour l'astuce. Cf. \xr{scalist}. --- \dv{oncophile} \textit{n.}

%EN ophelocracy
\E{ophélocratie}{ɔ.fe.lɔ.kʁa.si}{n. f.}{gr. ὄφελος profit, utilité + κράτος pouvoir} Gouvernement par l'utilité ; régime où l'on juge les institutions sur ce qu'elles livrent. --- \dv{ophélocrate} \textit{n.}

%EN opiferous
\E{opifère}{ɔ.pi.fɛʁ}{adj.}{lat. \la{opifer} qui porte secours} Qui apporte du secours ; se dit des techniques dont les bienfaits échoient surtout aux pauvres. --- \dv{opiférité} \textit{n. f.}

%EN opsitechny
\E{opsitechnie}{ɔp.si.tɛk.ni}{n. f.}{gr. ὀψέ tard + τέχνη art} L'arrivée tardive d'une technique qui eût pu venir des décennies plus tôt, n'eût été l'interdiction. \ex{L'opsitechnie est l'impôt silencieux.} \src{Tract. Cop.} vii.

%EN orbiculist
\E{orbiculiste}{ɔʁ.bi.ky.list}{n.}{lat. \la{orbiculus} petit cercle, d'après \textit{orbite}} Qui vit ou travaille en orbite. --- \dv{orbiculisme} \textit{n. m.} Cf. \xr{apoikiast}.

%EN orexergy
\E{orexergie}{ɔ.ʁɛk.sɛʁ.ʒi}{n. f.}{gr. ὄρεξις appétit + ἔργον ouvrage} L'appétit d'énergie ; la mesure de la disposition d'une société à consommer de la puissance dans la poursuite de ses fins. Cf. \xr{dynamophagy}.
